%
% "Simpulan dan Saran" dapat diganti menjadi "Simpulan dan Rekomendasi",
% "Simpulan", atau "Kesmpulan"
%
\section{Simpulan dan Saran}
%======================================================%
%        TULIS SIMPULAN DAN/ATAU SARAN DI SINI         %
%                                                      %
%          Mulailah menulis dari \subsection           %
%      \subsection > \subsubsection > \paragraph       %
%      > \subparagraph                                 %
%                                                      %
% Jangan Lupa untuk Menghapus Contoh Soal di Bawah Ini %
%======================================================%

Simpulan yang dituliskan pada bagian ini bukanlah sekadar ringkasan dari uraian yang telah dipaparkan sebelumnya. Narasi simpulan wajib bersumber dan disarikan secara padat dari temuan yang konkret serta penting, yang telah dibahas secara tuntas pada bagian hasil dan pembahasan. Penulis harus memperhatikan dua aturan krusial dalam menyusun bagian ini, yaitu tidak mengambil simpulan dari hal-hal yang tidak pernah dibahas pada bagian sebelumnya, serta tidak menyimpulkan sesuatu berdasarkan data yang kurang lengkap atau bahkan dari data yang tidak ada sama sekali. Seluruh argumentasi akhir disajikan secara lugas untuk menjawab perumusan masalah secara definitif.

Sementara itu, bagian saran atau rekomendasi dirumuskan untuk menyampaikan ide-ide baru yang timbul secara logis dari hasil pembahasan. Bagian ini memuat komentar serta masukan berharga yang ditujukan bagi peneliti lain atau sebagai bahan kajian lebih lanjut di masa depan. Rekomendasi disusun secara spesifik, realistis, dan memiliki keterkaitan erat dengan keterbatasan penelitian yang telah disebutkan di awal, sehingga memberikan arah yang jelas bagi pengembangan ilmu pengetahuan pada topik terkait selanjutnya.